% SAME-FAMILY WORKING FRAGMENT, 25 October 2026. Root-owned.
% General exact18 B16/D16 with ACTUAL typed defects. Outside review pending.
% Unintegrated/uncompiled; no scientific, Lean or TeX job.
\subsection{Typed defects beside \(B_{16}\) and \(D_{16}\)}
\label{app:BD16typed18}
These cases share the long-prefix and head geometry. They separate at
the final short layer: \(B_{16}\) has a genuine short normal and a
reciprocal endpoint, while \(D_{16}\) has two signed forks and a
centre calculation. We prove their common input once.

Use the actual unit roots \(n_i=(e_i-e_{i+1})/\sqrt2\) along
the long chain. For \(B_{16}\), \(n_{16}=e_{16}\),
\(\Omega_i=\sqrt2P_i\) for \(i\le15\), and
\(\Omega_{16}=P_{16}\). For \(D_{16}\),
\[
 n_{16}=(e_{15}+e_{16})/\sqrt2,\quad
 \Omega_i=\sqrt2P_i\ (i\le14),\quad
 \Omega_{15,16}=(P_{15}\mp e_{16})/\sqrt2.
\]
The last \(D\) coordinate is signed. Write the one original vector
as \(z=E(p-x-y)\), with the classical and rank-two projections
\(x,y\). The actual typed defect gives \(\norm y^2>.09\) and
\(\norm x^2+\norm y^2<1\).

Put \(c_7=\cos(\pi/7)\), \(s_7=c_7^2\).
The actual isolated defect wall \(d=A_Tp-v_T\) has
\[
\begin{array}{c|c|c}
&v_T&t_T=A_T^2\\\hline
N_7&c_7\omega_j^{A_2}&4s_7/3-1\\
O_7&\omega_j^{I_2(7)}/2&1/[4(1-s_7)]-1.
\end{array}
\]
The native-seven cubic gives \(.9009<c_7<.901\), hence
\(0<t_N<1/12\), \(.327<t_O<.342\).
Every classical opposite retains the whole old rank-two component.
An old \(I_2(7)\) accepts no contact: already a single three
has determinant \(1-s_7-1/4<0\). An old \(A_2\) accepts at
most one contact, with label three, four or five and whole cost
\(Q=1/3,2/3,\phi^2/3\). All these cases remain; no field
restriction removes the last one.

\subsubsection{Fresh tails, then strong five}
Set \(f_i=\sqrt i-.61/\sqrt2\). Appendix~\ref{app:prefixinventory}
reaches twelve in \(B_{16}\) and eleven in \(D_{16}\), leaving
only nonadjacent pure far-three helpers. We reuse the actual local
matrices and points in Appendix~\ref{app:exact18layers}, before
its ordinary-spectator consumer. Their proofs retain all spectators.
The positive foreign costs here are at least \(1/3\), so none
weakens a test that required only \(Q\ge1/4\).

In \(B\), the fresh \(A_{12}+B_3\) layer at thirteen has
exactly the long-three and short-three finite extensions printed
there. The latter is the sealed \(F_4\) white block; own zero
has its actual affine \(3,4,3,3\) path. The full long bridge
has residual \(1/26\); its minimum numerator, even at own zero,
is \(1/26+1/2\). No short-end bridge gives a finite long path.
Thus strong-or-pure-helper coverage reaches thirteen.

In \(D\), the fresh \(D_4\) layer at twelve retains its long
leaf, both distinct fork leaves, bridges and all foreign costs.
The only extra own-five/fork-three row has its actual finite
\(H_4\) point, flat through nine, with \(G>3.34,R<12\).
The local payment formula with \(K=8,9\) and the typed cost gives
\[
\begin{split}
 .7683275+.09+2(3.34-2.398)/7-4/49&>1.0458,\\
 .797425+.09+2(3.34-2.570)/7-3/49&>1.0457.
\end{split}
\]
The \(D_3\) layer at thirteen keeps middle-three and end-three,
four and five, including the middle/full-\(D_{16}\) bridge.
The own-five point has \(G>3.47,R<13\). For every remaining
finite own \(.9<g<1\), average the two actual \(A_2\) cores
whose corrections are \(g(4n_{13}/3+2n_{12}/3)\) and
\(g(4n_{13}/3+2n_{14}/3)\). Both satisfy every original wall.
With \(\alpha=\sqrt2(g+1/4)\), \(a=4g^2/3-1\), they give
\[
 A^2=13\alpha^2-5/8,\quad R<13,\quad
 G=(A-\sqrt a)/\alpha>13/4.
\]
Indeed its derivative has numerator
\((5/8)/A-(g/3+1)/\sqrt a<0\), and at one its value is
\(8\sqrt5/5-2\sqrt6/15>13/4\), using
\(\sqrt5>559/250,\sqrt6<49/20\).
Choosing actual strong \(K=9\) or \(10\) gives
\[
 .797425+.09+2(3.25-2.570)/7-4/49>1.000078,
\quad .823735+.09+2(3.25-2.732)/7-3/49>1.000510.
\]
This sharper comparison is needed for the typed spectator; the weaker
ordinary-case \(3.23\) estimate is not substituted. Coverage reaches
thirteen in both types. Helper packing extends because its actual
mixed-point proof uses the same inverse \(2\min(i,j)\) and positive
retained right blocks. It does not require those right blocks to have
the generic bulk ranks.

Some twelve/thirteen is strong. If five were weak, the actual
weak-five point gives \(x_1>.5965,S_6>2.01815\). The legal
interpolation through one, six and that strong end has energy at least
\[
 .7600299945+(\sqrt{12}-\sqrt6)^2/6
             =3.7600299945-2\sqrt2>.93.
\]
Adding only \(\norm y^2>.09\) is impossible. Thus five is strong,
before choosing a head or final centre. The ordinary references of
Appendix~\ref{app:exact18layers} now apply with every internal helper
omitted: \(q_{12}>.8689525,q_{13}>.88864\). Their first five
coordinates are \(b=f_5/5<.361\), with individual \(P_5\)
coefficient \(>.146\).

\subsubsection{A common head argument without arithmeticity}
Deleting long node one retains \(B_{15}\) or \(D_{15}\).
Finite classification allows only isolation or a simple three at
its long endpoint two. The latter has
\[
 c=Ap-g\Omega_1-\Omega_2/2-v_Z,\quad
 A^2=2g^2+2g+Q,\quad C_1=(2g+1)/\sqrt{1-2Q}.
\]
At own zero, zero foreign cost has zero time; positive foreign cost
has future null \(c+\Omega_2/2=\sqrt Qp-v_Z\).
The actual positive expansions of \(\Omega_2/2\) are
\(n_1+2(n_2+\cdots+n_{15})+\sqrt2n_{16}\) in \(B\), and
\(n_1+2(n_2+\cdots+n_{14})+n_{15}+n_{16}\) in \(D\).
Own three has its triangle. All other positive owns force
\(C_1\ge1+\sqrt2\) and \(x_1>.64\).
Graft the first five-block to \(.64,(5b-.64)/4,\ldots\).
Its other coordinates exceed \(.29\), above the next slope;
its \(P_5\) debit is less than \(.07<.146\). Its gain
\(>(5/4)(.279)^2>.0973\) makes
\(q_{12}+.0973+.09>1\). The head must be isolated within
the classical component, with foreign contacts still present.

Write \(c_H=Ap-g\Omega_1-v_Z\), \(A^2=2g^2+Q-1\).
For a finite pure \(N_7\) head its actual product with \(d\)
lies strictly between \(-1/2\) and zero. For a finite pure
\(O_7\) head its magnitude is less than \(\sqrt{.342}<1/\sqrt2\),
so it must equal \(1/2\). This forces
\[
 g^2=(1+1/(4t_O))/2,\quad148/171<g^2<577/654,
 \quad .93<g<.9394.
\]
But \(\cos(\pi/8)<.925\) and \(\cos(\pi/9)>.9395\),
the latter by \(8s^3-6s-1\). No finite native own lies there.

For a finite mixed \(N_7\) head the product is
\(q-A\sqrt{t_N}\), with
\(q=4\beta c_7/3\) or \(2\beta c_7/3\).
At \(\beta\ge1/\sqrt2\), \(q>.42\) but
\(A\sqrt{t_N}<\sqrt{2/12}<.409\). At \(\beta=1/2\)
the same-root value \(q>.6\) is also impossible. The other-root
product exceeds \(-1/2\) and must vanish. It gives
\(g^2=(3s_7-2)/(8s_7-6)=(1+1/(4t_O))/2\), the same gap.
Parallel has its actual future null. Every finite head is excluded.

For dotted mixed heads, \(Q\ge2/3\) forces \(C_1^2>6\)
and hence \(x_1>.64\). At \(Q=1/3\), same-root acuteness
gives \(A^2>108/25\) and \(C_1^2>187/25\).
In the other orientation
\[
 A^2t_N-q^2>(4/3)t_N-s_7/9=(5s_7-4)/3>0.
\]
The actual negative product therefore has magnitude at least \(1/2\),
so \(A\sqrt{t_N}>.8\), \(A^2>192/25\), \(C_1^2>313/25\).
Again \(x_1>.64\). Only a pure dotted head remains.

For it put \(u=g^{-2}\in(0,1)\). The actual positive core
\(c_H+g^{-1}n_1\) has square \(-(1-u)\), zero cap product,
negative own product, positive time and every omitted wall nonpositive.
It gives
\[
 \ip xZ\ge{\sqrt{2-u}-\sqrt{u(1-u)}\over\sqrt2},
 \quad Z=(1-u/2)e_1+(u/2)e_2.
\]
Pay weight \(\sqrt2/5\) on \(Z-P_5/5\), costing
\(\sqrt2/25<.057\) at the individual anchor. The residual
is zero-sum on the first five coordinates and orthogonal to the
reference. Its completed-square gain is
\[
 \eta=.4(\sqrt{2-u}-\sqrt{u(1-u)}-\sqrt2b)
                  -.04(1.6-2u+u^2)>.03624.
\]
For the full interval, \(\sqrt2b<.5106\). On the two halves use
lower root chords \(1.4142-.379u,1.4494-.4494u\) and upper
chords \(u^2\le u/2,1.5u-.5\). The lower gains become
\(.29744-.0916u-.4\sqrt{u(1-u)}\) and
\(.33152-.15976u-.4\sqrt{u(1-u)}\).
The circle inequality with radical upper bounds \(.4104,.4308\)
gives \(>.04644,>.03624\). This one paid head can be combined
with any later residual vanishing on the first five coordinates.

\subsubsection{The \(B_{16}\) last-long layer}
Delete fourteen, retaining the actual \(B_2\) on long fifteen
and short sixteen. A finite contact to it is a three at either end.
For long fifteen the Schur data are \(\sigma=g+1/2\),
\(\Delta=1/2-Q\). Own zero has null
\(c+n_{14}+2n_{15}+\sqrt2n_{16}\); own three has its triangle;
every other positive own is strong.
For short sixteen the data are \(\sigma=g+1/(2\sqrt2)\),
\(\Delta=1/2-Q\). Own zero has the actual affine \(F_4\)
kernel \((\sqrt2,2\sqrt2,3,2,1)\) on
\(c,n_{16},n_{15},n_{14},n_{13}\). Own four has null
\(\sqrt2c+n_{13}+2n_{14}+n_{15}\), six has
\(\sqrt3c+2n_{14}+n_{13}\), and parallel has \(c+n_{14}\).
Every remaining own at least five has ratio
\(\phi/\sqrt2+1/2>47/30>25/16\). Positive foreign
\(Q\ge1/3\) makes own three strong too.

All four left/right endpoint-three bridges are finite \(B_{16}\),
with residual \(1/28\) and \(Q=0\). Their minimum numerator,
including own zero and the short-right bridge, gives ratio
\[
 \sqrt{28}(1/28+1/(2\sqrt2))>27/14>25/16.
\]
Thus these own-zero bridges are directly estimated, not dismissed
by an invalid all-three cycle. Left-only and isolated cases use
their actual long-neighbour points through fifteen. The sole extra
weak nonhelper is consequently
\(c=Ap-(\Omega_{14}+\Omega_{16})/2\), with no foreign part.

For this wall put \(a=(1+\sqrt2)/2\). Its spatial part is
\(aP_{14}+(e_{15}+e_{16})/2\), and \(A^2=10+7\sqrt2\).
Release old fifteen instead. The retained chain
old one through fourteen, then \(c\), then the actual short sixteen
is finite \(A_{16}\), all bonds three. Positive finite projection
gives an actual future endpoint with the omitted old fifteen strict:
its products with the positive old-fourteen and short-sixteen terms
are negative. Every other wall is nonpositive termwise, and the full
rank-two block is retained. In the old edge direction
\(P_{15}/\sqrt{15}\) its data are
\[
 b_{15}={15+14\sqrt2\over2\sqrt{15}},\quad
 b_{15}^2-A^2=17/60,\quad
 C_{15}={15+14\sqrt2\over\sqrt{17}}>42/5.
\]
Therefore \(S_{15}>\sqrt{555/47}>3.43\) on the original vector.
Append this floor to an actual strong twelve/thirteen reference,
skipping intermediate unavailable floors. The slopes are
\(a_K=(3.43-f_K)/(15-K)\). For legality use
\(f_{12}>3.03,f_{13}>3.17\): then
\(a_{12}<.4/3<1/7\), \(a_{13}<.26/2<7/50\).
For energy use the opposite bounds
\(f_{12}<3.034,f_{13}<3.175\), gaining
\(>.396^2/3\) or \(>.255^2/2\).
The first five block and its head price are unchanged. One square
with the head and the one spectator gives totals
\(>1.0474645\) and \(>1.0473925\), respectively.
This disposes of the extra cap and extends helper coverage through
fourteen. Its actual mixed-point packing makes thirteen or fourteen
strong. Their energies exceed \(.88864\) and
\(.9084-4/3200=.90715\). The same head gain and spectator
then give at least \(.88864+.03624+.09>1\).
Thus the entire typed \(B_{16}\) row is impossible.

\subsubsection{The \(D_{16}\) centre and the point where cases split}
Delete centre fourteen, retaining \(A_{13}\), both actual fork
\(A_1\)'s and the full rank-two component. For own \(g\), left
contact \(\beta\) at \(j\), and fork contacts \(\delta,\epsilon\),
the exact data are
\[
 \Delta=1-Q-q_L-\delta^2-\epsilon^2>0,\quad
 \sigma=g+\beta j/14+(\delta+\epsilon)/2,\quad
 q_L=\beta^2j(14-j)/7,\quad C_{14}=\sqrt{28}\sigma/\sqrt\Delta.
\]
The sufficient strong ratio is less than \(25/16\).
The complete white list is: no left contact; left three at
\(1,2,12,13\) or left four at \(1,13\); left endpoint-three
plus one fork three/four or both forks three. Without a left
contact the two-fork pairs are \((3,3),(3,4),(3,5)\), up to
exchange; a single fork has any finite label. Full long bridges
are foreign-sealed. A permitted extra old-\(A_2\) arm is retained
when its full white Gram is positive.

The left-only and isolated roles have their fresh \(D_i=28\)
clearings. The isolated \(0<u<1\) loss is
\(<(.096+1.005)/2=.5505\). For \(Q>0,u>1\),
\(u\le8/(3\phi^2)<1.019\). Hence
\(\sqrt{28}-\sqrt{28-u}<.1\) and \(\sqrt{u-1}<.138\),
so the full loss is \(<.238\). These follow from
\(\phi>1.618,\sqrt{28}>5.29,\sqrt{28-1.019}>5.19\) and
\(.138^2>.019\). At \(u=1\) is the actual null. Pure five/seven/at-least-eight have respective
loss bounds \(25/160+5/24\), \(10/81+\sqrt{13/120}\),
and \(.511\). The left four, near-three and D-position tests
are stronger than \(25/16\). A positive foreign cost makes
far-left three strong by its actual full-foreign clearing. A pure
far-left helper at fourteen would make thirteen strong and give
\(q_{13}+.03624+.09>1\). No field argument enters these tests.

Left-three/fork-three bridges have residual \(2/7\), numerator
\(\ge g+2/7\); own zero/three have their cycles and owns at
least four are strong. Left-three/fork-four or two-fork bridges
have residual at most \(1/28\) and are directly strong including
own zero. Actual simple zero-own cycles, where present, are retained. No finite foreign arm fits.
For one fork-three, positive owns three/four/six/parallel have
the actual cycle, weighted-star, \(6,3\), or pair-null tests.
With \(Q\ge1/3\), all remaining owns are strong because
\[
 \Delta\le5/12,\quad \sigma\ge\phi/2+1/4,\quad
 \sigma^2/\Delta\ge3(8\phi+5)/20>625/256.
\]
Single fork at least four with own at least five is stronger.
Own-three/fork-four or five is retained below; foreign cost makes
it strong, and larger forks are directly strong. Own-zero/fork-three
has the affine \(E_8\) arms \(5,2,1\), fork-four has null
\(\sqrt2c+2n_f+2n_{14}+n_{13}+n_h\), and fork-six its
\(6,3\) null. Fork-five permits \(Q=0,1/3\); fork at least
seven forces \(Q=0\). For two forks own three has its triangle,
owns at least four are strong, and own-zero/\((3,3)\) is the
affine square. This leaves exactly the rows below.

Every remaining row gives a pure \(D\) pressure \(\ip xV\ge G\),
flat one on the first eleven coordinates. Choose actual strong
\(K=10\), or eleven if ten is weak. The terminal price \(1/7\)
and head price lie at different anchors; their residuals have
zero full cross because the centre residual vanishes on the
first \(K\) coordinates. The one-vector total is
\begin{equation}\label{eq:D16typedcentre}
 \norm x^2+\norm y^2>
 q_K+.03624+.09+2(G-f_K)/7-(R-K)/49.
\end{equation}

For positive own and single fork-three, \(Q=0\), put
\(\alpha=\sqrt2(g+1/4)\), \(A^2=14\alpha^2-3/4\).
Own five has its positive finite \(H_4\) point from
\(g(3n_{14}/2+n_{13}+n_{12}/2)\), giving \(G>3.60,R<14\).
For finite at least seven or dotted, weakness forces \(g<9/8\).
Use the actual core \(c+g(4n_{14}/3+2n_{13}/3)\), of square
\(-a\), \(a=4g^2/3-1\). Its old selected products vanish,
its cap product is \(-a\), both fork products are negative,
and every other wall is nonpositive termwise. Its displacement
has
\[
 R=14+(1/4-4g^2/3)/\alpha^2<14,\quad
 G=(A-\sqrt a)/\alpha>
 {2(\sqrt{835}-\sqrt{11})\over11\sqrt2}>3.28.
\]
The same derivative as above, now with \(3/4\), is negative,
so the endpoint \(9/8\) covers the entire weak interval.

For own zero and fork-five, \(Q=0\), project the actual \(H_4\)
on \(c,n_f,n_{14},n_{13}\). Its correction is
\(\delta(3n_f/2+n_{14}+n_{13}/2)\), \(\delta=\phi/2\).
Dividing its displacement by \(\delta/\sqrt2\) gives
\[
 V=P_{12}+(e_{13}+e_{14}+e_{15}\pm e_{16})/2,\quad
 R=13,\quad G>\sqrt{8\phi}-\sqrt{13-8\phi}>3.36.
\]
For the genuine \(N_7\) foreign-three row \(Q=1/3\), clear
the entire old \(A_2\) by its positive inverse combination \(v_Z\).
The finite point from \(c+v_Z+\delta n_f\), with
\(\Delta=2/3-\delta^2>0\), has displacement
\(\delta P_{14}/\sqrt2\). Since
\(A>2.137,\sqrt\Delta<.111,\delta<.81\), it gives
\(\ip x{P_{14}}>3.5\), so \(G>3.5,R=14\).
No foreign coordinate or second spectator energy is used.

For own zero and single fork at least seven, \(Q=0\), use
\(c+\delta^{-1}n_f+2(\delta^{-1}-\delta)n_{14}\).
Its cap/fork products vanish, its centre and other-fork products
are negative and its time is positive. With
\(z_0=\delta^{-2}\in(1,100/81)\) its data are
\[
 V=(1^{13},3-2z_0,z_0-1,\ \pm(1-z_0)),\quad
 G=\sqrt{16-2z_0}-\sqrt{-8+14z_0-6z_0^2},\quad
 R=24-16z_0+6z_0^2.
\]
The full interval has \(14G-R>32.51\). Indeed, for
\(z=z_0-7/6\), use
\(\sqrt{16-2z_0}>4.011-.27z_0\) and
\(\sqrt{1/6-6z^2}\le1/\sqrt6-3\sqrt6z^2\).
Endpoint squares and \(1.78^2<.036(96.8748)\) give
\(14G-R>32.528-1.78z+96.8748z^2>32.51\).

Own-three/fork-four or five has, after releasing that fork,
the actual positive finite \(D_{16}\) replacement and fork
cosine \(7+8\delta>12\). Own-zero/\((3,4)\) has the finite
\(B_{16}\) replacement, releasing the three-contacted fork,
with cosine \(7+4\sqrt2>12\). Positive inversion checks the
omitted fork as well as every other wall, so these are actual
endpoints. Both give
\(\ip x{P_{15}\pm e_{16}}>\sqrt{176/13}>3.67\), \(R=16\).
Finally own-zero/\((3,5)\) uses the actual finite \(H_3\)
projection on the cap and both forks. Its displacement is
\(\alpha P_{14}\), \(\alpha=(1+\phi)/(2\sqrt2)\), and
\[
 A^2=3+11\phi/2,\quad \Delta=(2-\phi)/4,\quad
 G={14+21\phi-\phi\sqrt7\over4A\alpha}>3.4,\quad R=14.
\]
The last strict comparison follows from
\(1.618<\phi<1.619,\sqrt7<2.646,A<3.451,\alpha<.927\):
its numerator exceeds \(43.696772\), whereas
\(3.4(4\cdot3.451\cdot.927)=43.5074472\).
Every finite projection uses its positive actual-root inverse,
has future time and retains all omitted original wall inequalities.

The smallest constant pressure, \(G>3.28,R<14\), makes
\eqref{eq:D16typedcentre} greater than \(1.0249\) at ten and
\(1.0248\) at eleven. The rows \((3.36,13),(3.5,14),
(3.67,16),(3.4,14)\) give larger margins. The continuous row
instead uses \((14G-R-14f_K+K)/49\), giving
\(>.949975+(32.51-38.248+10)/49>1.0369\), or
\(>.9734625+(32.51-40.404+11)/49>1.0368\).
Thus centre fourteen is strong. The ending-fourteen reference has
energy \(>.9084-4/3200=.90715\), retaining every permitted
internal hole. The same paid head and one typed spectator give
\(.90715+.03624+.09=1.03339>1\).
This excludes the typed \(D_{16}\) row as well. The entire
argument uses one original farthest vector, true short/fork normals,
individual prices and no arithmeticity or face inheritance.
